% --- LaTeX Essay Template - S. Venkatraman ---

% --- Set document class and font size ---

\documentclass[letterpaper, 11pt]{article}

% --- Package imports ---

% lipsum is just for generating placeholder text and can be removed
\usepackage{hyperref, lipsum} 

% --- Fonts (Set to Palatino or Times New Roman if desired) ---

% \usepackage{newpxtext, newpxmath}
% \usepackage{times}

% --- Page layout settiings ---

% Set page margins
\usepackage[left=1.2in, right=1.2in, top=.9in, bottom=.9in]{geometry}

% Anchor footnotes to the bottom of the page
\usepackage[bottom]{footmisc}
\usepackage{setspace}
\setstretch{1.15}
% Set line spacing
\renewcommand{\baselinestretch}{1.2}

% Use this for a one-off '*' footnote (e.g., a URL in a heading)
\newcommand{\starfootnote}[1]{%
\begingroup
\renewcommand{\thefootnote}{\fnsymbol{footnote}}%
\footnote{#1}%
\endgroup
\addtocounter{footnote}{-1}%
}

% Set spacing between paragraphs
\setlength{\parskip}{2mm}

% --- Page formatting settings ---



% --- Essay title and author ---
\title{\vspace{-1.5cm} Summary of \textit{Veridical Data Science}, Chapters 7 \footnote{\url{https://vdsbook.com/}}}
\author{Shizhe Zhang (3041882158)\\\href{mailto:shizhe\_zhang@berkeley.edu}{shizhe\_zhang@berkeley.edu}}

\date{\today}

% --- Main text ---
% --- Document starts here ---

\begin{document}
% Set link colors for labeled items and URLs (blue) 
\hypersetup{colorlinks=true, linkcolor=blue, urlcolor=blue}
\maketitle

\section*{Chapter 7: Clustering}

Chapter 7 introduces clustering as an unsupervised learning task aimed at discovering structure in data without labeled outcomes.
Unlike supervised methods, clustering does not have ground truth targets, and different algorithms and choices can yield different groupings of the same data.
Many of them can be equally valid depending on the context and goals of the analysis.


Then the chapter introduces different methods of clustering, such as K-means and hierarchical clustering,
and the metrics for evaluating cluster quality, such as silhouette scores and WSS.


K-means partitions data into K clusters by iteratively assigning points to the nearest cluster center and updating the centers,
But it depends on the choice of K, initialization.
Hierarchical clustering builds a tree structure that shows how data points can be grouped at different levels of dissimilarity.
Note that the measure of distance and the linkage method (e.g., single, complete, average) can lead to different clusterings.

We need to ensure the stability of clusters by checking how results change with different random seeds, preprocessing steps, and algorithm parameters.
A way to evaluate stability is to run the clustering multiple times and compare results using metrics that measure similarity between two clusterings.
The chapter discusses tools such as silhouette scores, WSS, stability analysis, and the Rand Index for comparing clustering results.

Overall, clustering is a powerful tool to help us understand the structure of data, but it requires careful consideration of the algorithm choice, parameter settings, and evaluation metrics.
% --- Document ends here ---
\end{document}
